\subsection{Design, Randomization, and Observational Comparisons}

Comparative inference is only as meaningful as the experiment or study that produced the groups.  The same two-sample statistic can estimate a causal treatment effect in a randomised experiment, describe an association in an observational study, or amplify bias if the groups were formed by a hidden mechanism.

\begin{definition} \label{def:controls-randomization-and-blinding} \hlt{Controls, Randomisation, and Blinding}
\begin{enumerate}[label=(\roman*)]
\setlength{\itemsep}{0pt}
\item A \hlt{controlled experiment} compares $Y\mid T=1$ and $Y\mid T=0$ under common conditions.
\item \hlt{Randomisation} assigns $T\in\{0,1\}$ so that $T$ is independent of pre-treatment factors $C$ by design.
\item \hlt{Blinding} conceals $T$ when knowledge of treatment could affect $Y$ or its measurement.
\end{enumerate}
\end{definition}

\begin{definition} \label{def:placebo-effect-and-confounding} \hlt{Placebo Effect and Confounding}
\begin{enumerate}[label=(\roman*)]
\setlength{\itemsep}{0pt}
\item A \hlt{placebo effect} is variation in $Y$ caused by belief in treatment rather than the active mechanism.
\item Placebos and blinding separate expectation effects from the causal effect of $T$.
\item A pre-treatment variable $C$ is \hlt{confounded} with treatment $T$ for response $Y$\\
if $\mathcal{L}(C\mid T=1)\neq\mathcal{L}(C\mid T=0)$ and $C$ affects $Y$.
\item An \hlt{observational study} compares groups without assigning $T$.\\
Group differences may reflect confounding rather than the factor of interest.
\end{enumerate}
\end{definition}

\begin{example} \label{ex:mammary-artery-ligation-design} \hlt{Mammary Artery Ligation}\\
A blinded trial randomised patients after preliminary incisions to ligation $(T=1)$ or sham surgery $(T=0)$. Similar outcomes show that controls remove background improvement; blinding limits placebo, observer effects.
\end{example}

\begin{principle} \label{prin:randomization-confounding-and-blocking} \hlt{Confounding, Blocking, and Pairing}\\
Randomisation balances unmeasured covariates $C$ in probability; blocking, matching, and pairing condition on known influential $C$ before randomisation. Pairing is one-factor blocking; blocked and factorial designs extend this control to multiple factors.
\end{principle}

\begin{figure}[H]
  \centering
  \begin{tikzpicture}[
  x=1cm,
  y=1cm,
  >=Stealth,
  line cap=round,
  line join=round,
  every node/.style={font=\scriptsize},
  box/.style={draw=black!70, rounded corners=2pt, align=center,
    minimum width=1.45cm, inner xsep=4pt, inner ysep=3pt},
  arrow/.style={->, black!72, line width=0.75pt},
  confound/.style={->, gray!65, dashed, line width=0.65pt},
  random/.style={->, blue!70!black, line width=0.9pt},
  block/.style={blue!70!black, line width=0.9pt},
  label/.style={text=black!70}
]
  \node[box] (factor) at (0,0) {$C$\\background};
  \node[box] (treat) at (2.9,0) {$T$\\treatment};
  \node[box] (resp) at (5.7,0) {$Y$\\response};
  \node[box, draw=blue!70!black, text=blue!70!black] (rand) at (2.9,-1.35) {randomize};

  \draw[arrow, bend left=22] (factor) to (resp);
  \draw[arrow] (treat) -- (resp);
  \draw[confound] (factor) -- (treat);
  \draw[random] (rand) -- (treat);

  \draw[block] (1.36,-0.17) -- (1.54,0.17);
  \draw[block] (1.54,-0.17) -- (1.36,0.17);
\end{tikzpicture}

  \caption{Randomization breaks the confounding path from $C$ to $T$.}
  \label{fig:randomization-confounding-paths}
\end{figure}

\begin{example} \label{ex:lanarkshire-and-cage-position-design} \hlt{Uncontrolled Adjustments and Position Effects}\\
In the Lanarkshire milk experiment, post-randomisation movement between groups made $T$ depend on perceived nourishment $C$, reintroducing selection bias; weighing children in clothes also confounded measurement with socioeconomic status. In the FD\&C Red No. 40 study, cage location was systematic in sex and dose, so position effects confounded treatment comparisons; randomising or blocking cage location would protect inference.
\end{example}

\begin{remark} \label{rem:comparative-study-checklist} \hlt{Comparative Study Checklist}\\
Before interpreting group comparisons, record sampling, treatment assignment $T$, controls, blinding, blocking or pairing factors $C$, dependence among repeated measurements, and the number of tested comparisons.\\
These determine whether the analysis supports causality, association, or hypothesis generation.
\end{remark}
